\documentclass[11pt,a4paper]{article}
\usepackage[utf8]{inputenc}
\usepackage[T1]{fontenc}
\usepackage{lmodern,microtype}
\usepackage{amsmath,amssymb,amsthm}
\usepackage{booktabs,array}
\usepackage[margin=2.4cm]{geometry}
\usepackage{xcolor}
\usepackage[colorlinks=true,linkcolor=blue!50!black,urlcolor=blue!50!black,citecolor=blue!50!black]{hyperref}
\usepackage{enumitem}
\setlist{itemsep=2pt,topsep=3pt}
\setlength{\parskip}{4pt}
\setlength{\parindent}{0pt}
\newcommand{\C}{\mathbb{C}}
\newcommand{\R}{\mathbb{R}}
\newcommand{\E}{\mathbb{E}}
\newcommand{\KG}{K_G^{\C}}
\newcommand{\ph}{\operatorname{phase}}
\newtheorem*{theorem*}{Theorem}
\newcommand{\lean}[1]{\texttt{\small #1}}

\title{\textbf{The complex Grothendieck constant}\\[2pt]
\large New bounds, machine-checked in Lean}
\author{Alin Galatan, with several AI research systems\\[2pt]
\small\url{https://github.com/agalatan/grothendieck-lean}}
\date{October 2026}

\begin{document}
\maketitle

\begin{center}
\fbox{\begin{minipage}{0.86\textwidth}\centering\vspace{4pt}
\Large $1.401691989174\;<\;\KG\;<\;1.404599424$\\[6pt]
\normalsize Both bounds are theorems of a Lean~4/Mathlib development, each resting on one explicitly stated
family of externally certified numerical inputs (Section~\ref{sec:trust}).\vspace{4pt}
\end{minipage}}
\end{center}

\begin{abstract}
Grothendieck's inequality bounds how much a semidefinite (vector) relaxation of a bilinear optimisation problem can exceed the
true optimum. The best constant in the complex case, $\KG$, has been known since 1987 to lie roughly in $[1.338,1.405]$; its
first two decimals were settled only in September 2026. We improve both bounds, to an interval of width $0.0029$, and verify
the arguments in Lean. This note explains the problem, the two constructions, exactly what the computer checks, and why the
remaining gap is hard: the two sides are dual optimisation problems, and the obstruction on both is the nonlinear (cubic and
higher) structure of bounded functions of Gaussian vectors.
\end{abstract}

\section{The question}
Let $A=(a_{ij})$ be an $m\times n$ complex matrix. Compare
\[
 C(A)=\sup_{|\varepsilon_i|=|\delta_j|=1}\Bigl|\sum_{i,j}a_{ij}\,\overline{\varepsilon_i}\,\delta_j\Bigr|,
 \qquad
 V(A)=\sup_{\|x_i\|=\|y_j\|=1}\Bigl|\sum_{i,j}a_{ij}\,\langle x_i,y_j\rangle\Bigr|,
\]
the first over unimodular scalars, the second over unit vectors in any complex Hilbert space. Clearly $C(A)\le V(A)$.
Grothendieck proved in 1953 that $V(A)\le K\,C(A)$ for a constant $K$ independent of $m$, $n$ and the dimension. The smallest
such $K$ is the \emph{complex Grothendieck constant} $\KG$.

\paragraph{Why it matters.}
\begin{itemize}
\item \emph{Optimisation.} $V(A)$ is a semidefinite program; $C(A)$ is a hard combinatorial problem. $\KG$ is the best
 approximation ratio of the semidefinite relaxation, and rounding algorithms for it are exactly upper-bound proofs.
\item \emph{Functional analysis.} $K_G$ is the norm of the identity between the injective and projective tensor
 norms on $\ell_\infty\otimes\ell_\infty$, and controls absolutely summing operators (Grothendieck's original setting).
\item \emph{Quantum information.} The real constant $K_G^{\R}$ is the largest ratio between quantum and classical values of
 two-player XOR games (Tsirelson); the finite-order constants $K_G(d)$ govern the nonlocality of qubit states.
\end{itemize}

\paragraph{History.}
\begin{center}\small
\begin{tabular}{@{}llll@{}}\toprule
year & who & lower bound & upper bound\\\midrule
1953 & Grothendieck & & $K<\infty$\\
1984 & Davie & $1.33807$ & \\
1987 & Haagerup & & $1.40490913\ldots$ (Gaussian phase rounding)\\
Sep 2026 & Guo, Fang, Lu & $1.35584631827$ & \\
Sep 2026 & Heilman, Jones, Malavolta & $1.4$ & $1.404898554746$\\
Oct 2026 & \textbf{this work} & $\mathbf{1.401691989174}$ & $\mathbf{1.404599424}$\\\bottomrule
\end{tabular}
\end{center}
Haagerup's integral
$1/\gamma^\dagger=1.404575934\ldots$, $\gamma^\dagger=\int_0^{\pi/2}\cos^2t\,(1+\sin^2t)^{-1/2}dt$, has been proposed as the
exact value. It lies inside our interval, $2.35\cdot10^{-5}$ below the upper bound. For the real constant the corresponding
recent results are $1.773\le K_G^{\R}\le1.7799$ (Heilman--Jones--Malavolta, October 2026).

\section{Two dual constructions}
Both bounds use the same Gaussian picture. Let $Z=(z_1,\dots,z_n)$ be a standard complex Gaussian vector and decompose
$L^2(\C^n,\gamma)$ into the \emph{Wick chaoses} $H_{p,q}$ (polynomials of bidegree $(p,q)$ in $z,\bar z$, orthogonalised).
Under the rotation $z\mapsto e^{i\theta}z$ every relevant object has \emph{charge one}, so only the bidegrees
$(1,0),(2,1),(3,2),\dots$ (degrees $1,3,5,\dots$) appear.

\subsection{Upper bounds: rounding schemes}
Given unit vectors, project a Gaussian onto them and round: $\varepsilon_i=f(\langle G,x_i\rangle)$ with $|f|\le1$. The
correlation of two rounded values depends only on $t=\langle x_i,y_j\rangle$,
\[
 \E\bigl[\overline{f(\cdot)}\,g(\cdot)\bigr]=H(t)=\sum_{k\ge0}c_k\,t|t|^{2k}.
\]
If $H(t)=c\,t+r(t)$ with $\|r\|_A=\sum_k|c_k-c\,\delta_{k0}|<c$, a tensor-power argument (Krivine's trick) gives
$\KG\le1/(c-\|r\|_A)$. Haagerup's choice $f=\ph$ gives $1.40491$. Our rounding is a smoothed two-coordinate phase rounding,
$\ph\bigl(z+\sqrt{1+a^2}\,v(|z|^2)\,w\bigr)$ with a twelve-term Laguerre profile $v$, multiplied by a damped polynomial phase
$e^{i\varphi}$; its kernel is $H(t)=\sum E[n,m]\,p(t)^n(-\tau(t))^m$ with an odd polynomial $p$ of degree $101$ and
$\tau(t)=(t^3+a^2t^7)/(1+a^2)$. Compared with $c\,t$, $c=625000/877873$, the residual is at most $133\cdot10^{-8}$:
\[
 \KG\le\frac1{c-133\cdot10^{-8}}=\frac{87787300000000}{62499883242891}<1.404599424 .
\]
The whole difficulty is to keep \emph{every} coefficient $c_k$, $k\ge1$, small at once: improving the cubic coefficient
$c_1$ tends to spoil $c_2,c_3,\dots$, and the infinite tail has to be bounded too.

\subsection{Lower bounds: Hermite projection games}
A lower bound needs an operator $T$ whose \emph{vector} value is close to $1$ while its \emph{scalar} value is small. We use
\[
 T=P_{1,0}-\epsilon\,P_{2,1},\qquad \epsilon=\tfrac{149}{200},
\]
the projection onto linear functions minus a multiple of the projection onto the cubic chaos.
\begin{itemize}
\item \emph{Vector value.} With $u_\ell=z_\ell/\|z\|$, $\sum_\ell\langle u_\ell,Tu_\ell\rangle\ge1-(1+\epsilon)/(n+1)\to1$.
\item \emph{Scalar value.} For all $|f|,|g|\le1$, $|\langle f,Tg\rangle|\le B=0.713423496548$.
\end{itemize}
A finite-partition transfer turns this into matrices, so $\KG\ge1/B>1.401691989174$.

The scalar cap is the hard part. Writing $h=\tfrac12(f+g)$, $k=\tfrac12(f-g)$,
\[
 \Re\langle f,Tg\rangle=t^2-\epsilon\|P_{2,1}h\|^2-\|P_{1,0}k\|^2+\epsilon\|P_{2,1}k\|^2,\qquad t=\|P_{1,0}h\|.
\]
Only the last term is positive: we must show that a \emph{bounded} function cannot have too large a cubic part without
paying elsewhere. After a rotation $P_{1,0}h=t\,z_1$, and the cubic part of $k$ splits into pieces that involve $z_1$
(\emph{visible} sectors) and a \emph{remote} residual $q\in H_{2,1}$ in the other coordinates. The proof then has three layers.
\begin{enumerate}
\item \textbf{Row reduction.} For each of 19 profiles $r$, the contribution of the visible sectors is the supremum of an
 explicit finite-dimensional Gaussian integral, a convex one-variable \emph{potential} $\Psi_r$, evaluated at the scale of
 the remote residual. One gate variable is shared across rows.
\item \textbf{The law of the remote cubic.} What remains is a statement about the distribution of $U=|q|^2/\E|q|^2$ for an
 arbitrary cubic Wick polynomial in arbitrarily many variables. We use exact identities and inequalities valid for
 \emph{every} such cubic: the fourth-moment cap $\E U^2\le62$; the sixth-moment bound
 $\E U^3\le9m-12+40.45313\,(m-2)^{3/2}$, $m=\E U^2$, which rests on thirteen tensor-contraction facts including the Wick
 expansion of the sixth cumulant; a first-moment bound; and families of \emph{Stein cuts}
 $\E\,G(U)\ge0$, obtained by Gaussian integration by parts and completing squares against spectral budgets, a covariance
 order and \emph{contraction-Gram budgets} (positive semidefinite Gram matrices of the coefficient tensor).
\item \textbf{A finite certificate.} The parameters are covered by $7{,}875$ boxes in $(t,\tau,x)$; on each box a positive
 combination of rows and node certificates ($336$ joint nodes, $9{,}237$ row nodes, each a polynomial majorant checked in
 exact rational arithmetic) bounds the scalar value by $B$.
\end{enumerate}

\subsection{The two problems are one}
For the real constant, Heilman, Jones and Malavolta proved an exact minimax: over rounding schemes $\mu$ with kernel
coefficients $c_m(\mu)$ and Hermite projection games $Q=\sum_m q_m\Pi_m$,
\[
 \frac1{K_G}=\sup_\mu\Bigl[c_1(\mu)-\sum_{m\ge3}|c_m(\mu)|\Bigr]=\inf_Q\ \mathrm{OPT}(Q),
\]
and the same duality holds in the complex case on charge-one bidegrees. An upper bound is a good point of the convex set of
achievable kernels; a lower bound is a supporting hyperplane of that set. There is \emph{no intrinsic gap}: the gap
between our bounds is entirely the gap between restricted constructions and the cost of certifying them.

\section{What exactly the computer checks}\label{sec:trust}
\begin{center}\small
\begin{tabular}{@{}>{\raggedright\arraybackslash}p{0.47\textwidth}>{\raggedright\arraybackslash}p{0.47\textwidth}@{}}\toprule
\textbf{Proved in Lean} & \textbf{Assumed (certified outside Lean)}\\\midrule
\multicolumn{2}{@{}p{0.96\textwidth}@{}}{\raggedright\emph{Lower bound.}\newline\lean{lower\_bound\_of\_nodeCaps : NodeCaps → IsGrothendieckBound K → 1.401691989174 < K}}\\[2pt]
the witness $T$ on $L^2$ of the Gaussian measure and its vector value; the transfer to matrices; the polarisation identity;
the per-row Gaussian reduction of $\langle f,Tg\rangle$ for every unimodular pair (frames, fibre decomposition, Parseval,
gates); all moment, Stein, spectral and contraction-Gram inequalities for every circular cubic, including the sixth-cumulant
Wick identity; the finite certificate (one \lean{native\_decide})
&
\lean{NodeCaps}: the 19 potentials $\Psi_r$ (defined in Lean) lie below their stored rational caps at the tabulated nodes.
These caps come from exact-arithmetic certificates (interval integration and anchored quadratic majorants), replayed
independently but not yet formalised.\\\midrule
\multicolumn{2}{@{}p{0.96\textwidth}@{}}{\raggedright\emph{Upper bound.}\newline\lean{complex\_grothendieck\_upper : KernelHyps E EX E0 → IsGrothendieckBound (87787300000000/62499883242891)}}\\[2pt]
the transfer inequality (tensor powers); kernel composition and locality; the head of the kernel through degree 501 from
21,168 stored energies by a verified interval evaluator (one \lean{native\_decide}); the tail bound; the endpoint
&
\lean{KernelHyps}: the stored tables are the energies of the rounding (validity, summability, parity); the comparison
between the rounding and its truncation ($93.9\%$ of the residual budget); a weighted square-sum bound for the tail
($5.4\%$); two table budgets. The part computed in Lean is $0.54\%$.\\\bottomrule
\end{tabular}
\end{center}
Axioms: \lean{propext}, \lean{Classical.choice}, \lean{Quot.sound}, and exactly two pinned \lean{native\_decide} facts. No
\lean{sorry}; CI builds everything and checks the axiom report (140 theorems). The development has about $28{,}000$ lines of
hand-written Lean and $11{,}000$ lines of generated certificate data. The lower bound is close to fully formal: only
one-variable numerical caps remain. The upper bound still rests mainly on an external analytic comparison lemma.

\section{Why the gap is hard}
\paragraph{The witness class caps the lower bound.} A game supported on degrees $\{1,3\}$ sees only the cubic constraint on
roundings. Even a perfect evaluation of the norm of $P_{1,0}-\epsilon P_{2,1}$ cannot certify more than about $1.40226$
(numerical, best $\epsilon\approx0.756$). Proved obstructions are known for higher degrees: every game of degree $\le5$ is
capped at $1.4036705$, and every game of degree $\le7$ at $1.4044895$, both below Haagerup's value. Reaching the upper
end of the interval therefore needs games of degree $\ge9$; in the real case the certified game uses all odd degrees up to
$21$.

\paragraph{Certifying higher degree is the real bottleneck.} Two different certificate architectures, gated rows and
support functions, both lose about $2\cdot10^{-2}$ on the best degree-5 game. The reason is structural: a bounded function
with a large cubic part is forced to have a degree-five part too. Through the visible coordinate, the cubic part sees the
average of the remote coefficient profile and the quintic part its variation, and boundedness makes the profile vary. A
certificate that treats the chaos degrees as independent cannot see this coupling.

\paragraph{Why the complex case is harder than the real one.} For real $\pm1$ responses, $h=\frac12(f+g)$ and
$k=\frac12(f-g)$ have disjoint supports, so the two sides of a game decouple at every degree. For complex unimodular
responses $|h|^2+|k|^2=1$ and $\Re(h\bar k)=0$ force $f=e^{i\alpha}e^{i\theta}$, $g=e^{i\alpha}e^{-i\theta}$: one shared
phase field $\alpha$ drives both sides, and they never separate. The same shared phase is why marginal laws do not
determine the relevant nonlinear correlations: $(H(X),X)$ and $(H(Y),X)$, with $H(X)=X(|X|^2-2)/\sqrt2$ and $X,Y$
independent, have identical marginals and covariances, but $\E[\bar r\,q|q|^2]$ equals $8\sqrt2$ in the first case and $0$
in the second.

\section{Open problems}
\begin{enumerate}
\item Formalise the node caps (\lean{NodeCaps}), which would make the lower bound fully machine-checked.
\item Formalise the rounding behind the upper bound: its Hermite expansion, the comparison lemma and the contour tail bound.
\item Certify a game of degree $\ge5$ with a \emph{multi-grade fibre}: one bounded function of the visible coordinate and one
 remote direction per degree, so that the cross-degree coupling is kept. This is the route to a lower bound above $1.403$.
\item Run a complex primal/dual search, alternating linear programs over games of degree $\le D$ and pools of roundings, to
 locate $\KG$ to about $10^{-4}$ and to test whether Haagerup's value $1/\gamma^\dagger$ is the answer.
\item Prove the sharp first-moment bound $\E|q|\le\frac{\sqrt\pi}{2}\|q\|_2$ for cubic Wick polynomials. It is natural, it
 appears to be open, and it is the largest single lever on the law side.
\end{enumerate}

\section*{How this was made}
The constructions and certificates were found by several AI research systems working on a shared manuscript, cross-checking
each other's packages by independent replay. This Lean development turns the certificates into theorems: every claim above
marked as proved is a Lean theorem whose axioms are checked in CI. The research manuscript and the certificate packages are
kept outside the repository.

\section*{References}
\small\sloppy
\begin{itemize}
\item A.~Grothendieck, \emph{R\'esum\'e de la th\'eorie m\'etrique des produits tensoriels topologiques}, 1953.
\item U.~Haagerup, \emph{A new upper bound for the complex Grothendieck constant}, Israel J.\ Math.\ 60 (1987).
\item A.~M.~Davie, \emph{Lower bound for $K_G$}, unpublished note, 1984.
\item S.~Heilman, C.~Jones, G.~Malavolta, \emph{Sharper bounds for the complex Grothendieck constant},
 \href{https://arxiv.org/abs/2609.25462}{arXiv:2609.25462}.
\item S.~Heilman, C.~Jones, G.~Malavolta, \emph{Primal/dual method for the Grothendieck constant},
 \href{https://arxiv.org/abs/2610.10477}{arXiv:2610.10477}.
\item Guo, Fang, Lu, \href{https://arxiv.org/abs/2609.07000}{arXiv:2609.07000}.
\item R.~Saha et al., \emph{New lower and upper bounds for the Grothendieck constant},
 \href{https://arxiv.org/abs/2608.11158}{arXiv:2608.11158}.
\item A.~Naor, O.~Regev, \emph{Krivine schemes are optimal}, Proc.\ AMS 142 (2014),
 \href{https://arxiv.org/abs/1205.6415}{arXiv:1205.6415}.
\item G.~Pisier, \emph{Grothendieck's theorem, past and present}, Bull.\ AMS 49 (2012),
 \href{https://arxiv.org/abs/1101.4195}{arXiv:1101.4195}.
\item N.~Brunner et al., \emph{Bell nonlocality}, Rev.\ Mod.\ Phys.\ 86 (2014),
 \href{https://arxiv.org/abs/1303.2849}{arXiv:1303.2849}.
\end{itemize}
\end{document}
